% GENERATED FILE. Edit canonical lessons, then run npm run generate:books.
% canonical-source-hash: fnv1a64:79fb22c8ea9b0b88
% canonical-lessons: SA-C246-lajah, SA-C246-dhanyam, SA-C246-saktuh, SA-C246-ksaudram, SA-C246-mistannam

\chapter{Puffed rice, Grain, Barley flour, Wild honey, Sweet dish}
\label{ch:sa-puffed-rice-grain-barley-flour-wild-honey-sweet-dish}
\hlchaptermodality{246}

\begin{chapteropening}
\textbf{By the end of this chapter:} \emph{I can say puffed rice, grain, barley flour, wild honey and a sweet dish.}

Complete the last lesson of chapter 246: I can say puffed rice, grain, barley flour, wild honey and a sweet dish.
\end{chapteropening}

\section[\texorpdfstring{\sk{लाजाः} (lājāḥ)}{लाजाः (lājāḥ)}]{\sk{लाजाः} (lājāḥ) — puffed rice}

\label{lesson:SA-C246-lajah}



\begin{quote}
Before the new one: say the Sanskrit for a clove, then the Sanskrit for cardamom.
\end{quote}

\subsection*{You'll want to know: \sk{लाजाः}}
\textbf{\sk{लाजाः}} — \emph{lājāḥ} — \textquotedblleft{}puffed rice\textquotedblright{}.

\begin{grammarlens}[title={What you've built}]
One more word for this chapter.
\end{grammarlens}

\subsection*{Your turn}
\begin{itemize}
\raggedright
  \item \emph{Say it:} \emph{lājāḥ}
  \item \emph{Say it:} \emph{lājāḥ}, once more
  \item \emph{Say it:} say \emph{lājāḥ}
  \item \emph{Recall:} say the Sanskrit for a clove, then the Sanskrit for cardamom, then say \emph{lājāḥ} again
\end{itemize}

\begin{tcolorbox}[breakable,colback=teal!4,colframe=teal!35!black,title={Before you move on}]
What does \sk{लाजाः} mean? (\textquotedblleft{}Puffed rice\textquotedblright{}.) Say it once more.
\end{tcolorbox}
\section[\texorpdfstring{\sk{धान्यम्} (dhānyam)}{धान्यम् (dhānyam)}]{\sk{धान्यम्} (dhānyam) — grain}

\label{lesson:SA-C246-dhanyam}



\begin{quote}
Before the new one: say the Sanskrit for cardamom, then the Sanskrit for puffed rice.
\end{quote}

\subsection*{You'll want to know: \sk{धान्यम्}}
\textbf{\sk{धान्यम्}} — \emph{dhānyam} — \textquotedblleft{}grain\textquotedblright{}.

\begin{grammarlens}[title={What you've built}]
One more word for this chapter.
\end{grammarlens}

\subsection*{Your turn}
\begin{itemize}
\raggedright
  \item \emph{Say it:} \emph{dhānyam}
  \item \emph{Say it:} \emph{dhānyam}, once more
  \item \emph{Say it:} say \emph{dhānyam}
  \item \emph{Recall:} say the Sanskrit for cardamom, then the Sanskrit for puffed rice, then say \emph{dhānyam} again
\end{itemize}

\begin{tcolorbox}[breakable,colback=teal!4,colframe=teal!35!black,title={Before you move on}]
What does \sk{धान्यम्} mean? (\textquotedblleft{}Grain\textquotedblright{}.) Say it once more.
\end{tcolorbox}
\section[\texorpdfstring{\sk{सक्तुः} (saktuḥ)}{सक्तुः (saktuḥ)}]{\sk{सक्तुः} (saktuḥ) — barley flour}

\label{lesson:SA-C246-saktuh}



\begin{quote}
Before the new one: say the Sanskrit for puffed rice, then the Sanskrit for grain.
\end{quote}

\subsection*{You'll want to know: \sk{सक्तुः}}
\textbf{\sk{सक्तुः}} — \emph{saktuḥ} — \textquotedblleft{}barley flour\textquotedblright{}.

\begin{grammarlens}[title={What you've built}]
One more word for this chapter.
\end{grammarlens}

\subsection*{Your turn}
\begin{itemize}
\raggedright
  \item \emph{Say it:} \emph{saktuḥ}
  \item \emph{Say it:} \emph{saktuḥ}, once more
  \item \emph{Say it:} say \emph{saktuḥ}
  \item \emph{Recall:} say the Sanskrit for puffed rice, then the Sanskrit for grain, then say \emph{saktuḥ} again
\end{itemize}

\begin{tcolorbox}[breakable,colback=teal!4,colframe=teal!35!black,title={Before you move on}]
What does \sk{सक्तुः} mean? (\textquotedblleft{}Barley flour\textquotedblright{}.) Say it once more.
\end{tcolorbox}
\section[\texorpdfstring{\sk{क्षौद्रम्} (kṣaudram)}{क्षौद्रम् (kṣaudram)}]{\sk{क्षौद्रम्} (kṣaudram) — wild honey}

\label{lesson:SA-C246-ksaudram}



\begin{quote}
Before the new one: say the Sanskrit for grain, then the Sanskrit for barley flour.
\end{quote}

\subsection*{You'll want to know: \sk{क्षौद्रम्}}
\textbf{\sk{क्षौद्रम्}} — \emph{kṣaudram} — \textquotedblleft{}wild honey\textquotedblright{}.

\begin{grammarlens}[title={What you've built}]
One more word for this chapter.
\end{grammarlens}

\subsection*{Your turn}
\begin{itemize}
\raggedright
  \item \emph{Say it:} \emph{kṣaudram}
  \item \emph{Say it:} \emph{kṣaudram}, once more
  \item \emph{Say it:} say \emph{kṣaudram}
  \item \emph{Recall:} say the Sanskrit for grain, then the Sanskrit for barley flour, then say \emph{kṣaudram} again
\end{itemize}

\begin{tcolorbox}[breakable,colback=teal!4,colframe=teal!35!black,title={Before you move on}]
What does \sk{क्षौद्रम्} mean? (\textquotedblleft{}Wild honey\textquotedblright{}.) Say it once more.
\end{tcolorbox}
\section[\texorpdfstring{\sk{मिष्टान्नम्} (miṣṭānnam)}{मिष्टान्नम् (miṣṭānnam)}]{\sk{मिष्टान्नम्} (miṣṭānnam) — a sweet dish}

\label{lesson:SA-C246-mistannam}



\begin{quote}
Before the new one: say the Sanskrit for barley flour, then the Sanskrit for wild honey.
\end{quote}

\subsection*{You'll want to know: \sk{मिष्टान्नम्}}
\textbf{\sk{मिष्टान्नम्}} — \emph{miṣṭānnam} — \textquotedblleft{}a sweet dish\textquotedblright{}.

\begin{grammarlens}[title={What you've built}]
One more word for this chapter.
\end{grammarlens}

\subsection*{Your turn}
\begin{itemize}
\raggedright
  \item \emph{Say it:} \emph{miṣṭānnam}
  \item \emph{Say it:} \emph{miṣṭānnam}, once more
  \item \emph{Say it:} say \emph{miṣṭānnam}
  \item \emph{Recall:} say the Sanskrit for barley flour, then the Sanskrit for wild honey, then say \emph{miṣṭānnam} again
\end{itemize}

\begin{tcolorbox}[breakable,colback=teal!4,colframe=teal!35!black,title={Before you move on}]
What does \sk{मिष्टान्नम्} mean? (\textquotedblleft{}A sweet dish\textquotedblright{}.) Say it once more.
\end{tcolorbox}
